\subsection{Semi-infinite construction}

\begin{Definition} \label{eq: sec3 eval def}The evaluation representation $V_{u}$ of the algebra $\qD$ is a vector space with the basis $x^{k}$ for $k \in \mathbb{Z}$ and the action
\begin{gather*}
E_{a,b} x^{k} = u^a q^{\frac{ab}{2} + ak} x^{k+b}, \qquad c=c'=0.
\end{gather*}
\end{Definition}

\begin{Remark}The associative algebra $\qaD$ acts on $V_u$. The representation of $\qD$ is obtained via evaluation homomorphism $\mathrm{ev}\colon \qD \rightarrow \qaD$.
\end{Remark}

\begin{Remark}Informally, one can consider $x^k \in V_u$ as $x^{k-\alpha}$ for $u=q^{-\alpha}$. Define the action of $\qaD$ as follows. The generator $x$ acts by multiplication and $D x^{k-\alpha} = q^{k-\alpha} x^{k-\alpha} = u q^k x^{k-\alpha}$. However, $q^{-\alpha}$ is not well defined for arbitrary complex $\alpha$. So we consider $u$ as a parameter of representation instead of $\alpha$.
\end{Remark}

Let us consider the \emph{semi-infinite exterior power} of the evaluation representation $\siL V_{u}$. It is spanned by $| \lambda,l \rangle = x^{l-\lambda_1} \wedge x^{l+1-\lambda_2} \wedge \cdots \wedge x^{l+N} \wedge x^{l+N+1} \wedge x^{l+N+2} \wedge \cdots$, where $\lambda$ is a Young diagram and $l \in \mathbb{Z}$. Let $p_1 > \dots> p_i$ and $q_1> \dots> q_i$ be Frobenius coordinates of $\lambda$.

\begin{Proposition} \label{Prop:SemiInf-Fermi}
There is a $\qD$-modules isomorphism $\siL V_{u} \xrightarrow{
\,\smash{\raisebox{-0.65ex}{\ensuremath{\scriptstyle\sim}}}\,} \mathcal{M}_{u}$ given by
\begin{gather}\label{eq:lambda l}
| \lambda,l \rangle \mapsto (-1)^{\sum_k (q_k-1)} \psi_{-p_1+l} \cdots \psi_{-p_i+l} \, \psi_{-q_i-l+1}^* \cdots \psi_{-q_1-l+1}^* \n.
\end{gather}
\end{Proposition}

\begin{Proposition} \label{Boson-Fermion}
There is an isomorphism of $\qD$-modules $\mathcal{M}_u \cong \bigoplus_{l \in \mathbb{Z}} \mathcal{F}_{q^l u}$.
The sub\-module~$\mathcal{F}_{q^l u}$ is spanned by $| \lambda,l \rangle$.
\end{Proposition}

\begin{proof}Recall the ordinary boson-fermion correspondence (see \cite{KR}). The coefficients of
\begin{gather*}
a(z) = \sum_{n} a_n z^{-n-1} = : \! \psi(z) \psi^* (z) \! :_{(0)}
\end{gather*}
are indeed generators of the Heisenberg algebra. Moreover, $F^{\psi} = \oplus_{l \in \mathbb{Z}} F^{a}_{-l}$. The highest vector of $F^{a}_{-l}$ is $\n$ (in particular, $a_0 \n = -l \n $). Note that this is the decomposition of $\qD$-modules as well. Also, note that
\[\langle l | E(z) | l \rangle = \frac{q^{l} u }{1-q}, \qquad \langle l | F(z) | l \rangle = \frac{q^{-l} u^{-1} }{1-q^{-1}}.\]
Therefore one can use Proposition~\ref{Fock} for each summand $F^{a}_{-l}$.
\end{proof}

